\section{From task substitution to a workforce hypothesis}
\label{sec:implications}

The observations support a positive, bounded conclusion: software can execute the specified
governance-review operations in this engineered environment. The deterministic control makes
the case for an executable decision kernel especially clear. The agents demonstrate another
implementation of many of these operations, with measurable differences in evidence handling,
reliability, and cost. Neither result is weakened by the fact that software uses several kinds
of components. What remains unresolved is the effort needed to construct and operate the
interfaces when facts and policies arrive in less prepared forms.

This is also where the claim about human replacement becomes testable. If a deployed system
performs accepted tasks previously executed by people and satisfies the labor inequality at
unchanged output and quality, it has substituted for part of their work. Calling the remaining
people supervisors does not restore the eliminated execution hours. At fixed demand, fewer
required hours can support fewer posts; actual headcount also depends on allocation, utilization,
redeployment, and new work. The argument therefore permits substantial workforce displacement
without treating it as a measured consequence of this synthetic experiment.

\subsection{The 2033 hypothesis and a falsifiable deployment test}

The author's original hypothesis is that the DGF ecosystem will require at least 80\% fewer
workload-equivalent FTE by 20 September 2033 than during the reference year ending on
20 September 2026, at comparable governed output, quality, and service. The boundary includes
preparation, exceptions, verification, remediation of review errors, maintenance, supervision,
and supplier support. On an illustrative 140-FTE baseline, the target is at most 28 FTE.
This is an author-selected forecast, not an extrapolation fitted to benchmark scores.

The baseline window has already closed and no enterprise cohort has been enrolled. Testing
this exact dated hypothesis requires auditable historical activity records and a declared
population and aggregation rule before observing its endpoint. Otherwise it cannot be scored
as confirmed. A newly registered baseline can support a separate seven-year prospective test,
but cannot move the date of the original prediction. A credible test fixes the accounting
boundary, useful-hours conversion, output mix, quality tolerances, and service levels. It
reports the complete $L_A/L_H$ account rather than headcount in a renamed department. A ratio
above 0.20 at the deadline fails the proposed 80\% target; achieving that ratio while degrading
the fixed quality or service requirements also fails the substitution claim. Missing records
make the result unassessable, not successful.

Full elimination of human DGF execution is a stronger, undated conjecture. It would require
both case-level human effort and the recurring support floor to disappear. No finite benchmark
here establishes that endpoint. The scientifically actionable question is how much accepted
task execution can already be transferred and what prevents further transfer.

\subsection{The next discriminating experiment}

The most useful extension varies information preparation while preserving answerability.
Compare agents receiving structured snapshots with agents receiving the source documents and
operational records from which the same facts can be recovered. Separately vary executable
versus natural-language policy. Fix source access, mandate, time and tool budgets, and the
evidence contract before testing on held-out cases. Include the rule engine and any required
extraction pipeline in the comparison. Report decision quality, supported findings, complete
routes, abstention, and all engineering and operating costs.

This extension does not require claiming human-level accuracy. It asks which parts of a review
the system can perform when the required facts are no longer supplied as a ready-to-use snapshot.
Our counterexample means that simply deleting snapshots is not a valid design. The released
source audit also identifies runtime tools that can return snapshots, so changing file visibility
alone would not isolate the intended condition. The existing extraction ablation and General-only
intervention are preparations, not additional measured results in this paper.

\subsection{Limits of the present evidence}

The benchmark's policy, generator, and reference share an author-designed rule system; their
agreement is internal consistency rather than validation of a universal business standard.
Decision-balanced synthetic cases do not estimate enterprise prevalence. There are three
endpoints, one original trajectory per model and case, and repetitions on only 15 existing
dossiers. The post-hoc structural evidence metric is not semantic adjudication. No measured
human comparator, enterprise deployment, staffing effect, or comparative ranking of candidate
business functions is available. These limits locate the evidence: execution of declared review
contracts is measured; net labor substitution and the 2033 forecast remain hypotheses to test.
